\documentclass{leveldoc}
\DocUses{foundations,first_bend}
\DocFingerprint{65f43923b1ceeac5}

\begin{document}

\section{What the level asks}\label{sec:asks}

The game's text:
\leveldescription

The particle is launched straight up from $(6, 3)$, but the detector lies behind it, in the
corner below and to the left. A fixed charge $Q_1 = \val{QOne}$ at $\val{rOne}$ attracts it,
and not enough to turn it around. The player places up to three charges (as in
\levelref{first_bend}) to bring the particle back into the detector before it touches a
charge, leaves the world or runs out of time (\xsecref{foundations}{sec:board}).

The reference solution places $Q_2 = \val{QTwo}$ at $\val{rTwo}$: together the two attractions
bind the particle. It climbs, slows almost to rest, and falls back into the detector through its
top face. With two centres the orbit has no closed form, so this document computes the flight
numerically, with a method a person could carry out by hand, and with an error estimate.

\section{Setup}\label{sec:setup}

The level file, field by field (the appendix reproduces it):
\setuptable

In symbols: the particle has charge $q = \val{q}$, mass $m = \val{m}$, radius $a = 0$; the
speed of light is $c = \val{c}$. It starts at $\vect{x}_0 = (6, 3)$ with kinetic energy
$T_0 = \val{Tzero}$ along the direction of the level file, $(\val{dirX}, 1)$: the computed
$(\cos 90^\circ, \sin 90^\circ)$, whose first component, a rounding error, moves the flight by
less than $10^{-14}$ cells. The charges are $Q_1 = \val{QOne}$ at $\vect{r}_1 = \val{rOne}$ (the
level's) and $Q_2 = \val{QTwo}$ at $\vect{r}_2 = \val{rTwo}$ (the reference's), each of radius
$R = \val{R}$, with couplings $\kappa_i = qQ_i$: $\kappa_1 = \val{kOne}$ and
$\kappa_2 = \val{kTwo}$ (both attractive). The detector is the rectangle $0 \le x \le 3$,
$0 \le y \le 4$; the world is $-1 \le x \le 31$, $-1 \le y \le 21$; $t_{\max} = \val{tmax}$.

\begin{figure}[htbp]
\centering
\begin{tikzpicture}
\begin{axis}[axis equal image, width=0.92\linewidth, xmin=-1.6, xmax=31.6, ymin=-1.6, ymax=21.6,
  xtick={0,5,...,30}, ytick={0,5,...,20}, tick label style={font=\small},
  xlabel={$x$}, ylabel={$y$}, label style={font=\small}]
  \draw[black!8] (0,0) grid (30,20);
  \draw[black!45, dashed] (-1,-1) rectangle (31,21);
  \fill[green!50!black, fill opacity=0.18] (0,0) rectangle (3,4);
  \draw[green!40!black] (0,0) rectangle (3,4);
  \addplot[thick, dashed, orange!80!black] table {zvc.dat};
  \addplot[very thick, blue!70!black] table {path.dat};
  \fill[blue!75!black] (15,12) circle[radius=0.3];
  \node[font=\small, blue!60!black, anchor=south west] at (15.3,12.2) {$Q_1$};
  \fill[blue!75!black] (2,0) circle[radius=0.3];
  \node[font=\small, blue!60!black, anchor=south west] at (2.3,0.2) {$Q_2$};
  \draw[-{Stealth[length=6pt]}, thick] (6,3) -- (6,5.5);
  \fill (6,3) circle[radius=0.15];
  \node[font=\small, anchor=west] at (6.3,3.3) {launch};
  \node[font=\small, green!30!black, anchor=north] at (1.5,-0.1) {detector};
\end{axis}
\end{tikzpicture}
\caption{The board with the reference solution: the two attracting charges, the computed flight
(§\ref{sec:hand}) from the launch to the entry into the detector, and the edge of the region the
particle's energy allows (dashed, \eqref{eq:allowed}). The flight turns where it nearly touches
that edge.}
\label{fig:board}
\end{figure}

\section{Principles and model}\label{sec:model}

\begin{principles}
  \item The fields of charged spheres, superposed, and the force on a non-polarizable particle,
    \xeqref{foundations}{eq:coulomb} and \xeqref{foundations}{eq:sphereforce}.
  \item The relativistic equation of motion \xeqref{foundations}{eq:eom} and the launch
    \xeqref{foundations}{eq:launch}; the conservation of energy \xeqref{foundations}{eq:energy}.
  \item The board's rules (\xsecref{foundations}{sec:board}).
  \item The classical Runge--Kutta method, cubic Hermite interpolation and the error estimate
    by halving the step (\xsecref{foundations}{sec:numerics}).
\end{principles}

With the charges $Q_i$ at $\vect{r}_i$, the state $(\vect{x}, \vect{p})$ obeys
\begin{equation}\label{eq:eom-two}
  \deriv{\vect{x}}{t} = \frac{\vect{p}}{\gamma m}, \qquad
  \deriv{\vect{p}}{t} = \sum_{i=1}^{2} \kappa_i\,\frac{\vect{x}-\vect{r}_i}{\abs{\vect{x}-\vect{r}_i}^3},
  \qquad \gamma = \sqrt{1 + \frac{\abs{\vect{p}}^2}{m^2c^2}},
\end{equation}
from $\vect{x}_0$ and $\vect{p}_0 = \abs{\vect{p}_0}(\val{dirX}, 1)$, until the first event. The
motion stays in the plane (\xsecref{foundations}{sec:plane}). The assumptions are those of
\levelref{first_bend}, for two charges:

\begin{assumption}{The charges are held in place}\label{as:fixed}
Neither charge recoils (a rule of the game).
\end{assumption}

\begin{assumption}{No radiation}\label{as:radiation}
The radiated energy \xeqref{foundations}{eq:larmor} is not taken from the particle. Along the
reference flight it is $\val{Erad}$, a fraction $\val{EradFrac}$ of the launch energy
(§\ref{sec:hand}): its effect on the flight is of that relative order, far below the accuracy of
the game's integration or of this document's.
\end{assumption}

\begin{assumption}{A point particle}\label{as:point}
$a = 0$, and the particle does not polarize: \xeqref{foundations}{eq:sphereforce} is exact.
\end{assumption}

\section{Derivation}\label{sec:derivation}

\subsection{Why there is no orbit equation}

With one charge, the conservation of angular momentum turned the flight into a single linear
equation for $u = 1/r$, \xeqref{first_bend}{eq:orbit}. With two the force is not central, and
the angular momentum about either charge changes. In Newton's mechanics the problem of two fixed
centres can still be solved (Euler's problem): in elliptic coordinates
$(\sigma,\tau)$ with the charges at the foci, $r_{1,2} = d(\sigma \pm \tau)$ ($2d$ apart), the
Hamilton--Jacobi equation separates (\Goldstein{Ch.~10} for the method). Relativistically, the
square of the energy relation brings in $V^2$, and the product
$(\sigma^2-\tau^2)\,V^2$ that the separation needs to split into a function of $\sigma$ plus one
of $\tau$ contains
\begin{equation}\label{eq:noseparation}
  \frac{\kappa_1^2}{d^2}\,\frac{\sigma - \tau}{\sigma + \tau} + \frac{\kappa_2^2}{d^2}\,\frac{\sigma + \tau}{\sigma - \tau},
\end{equation}
which does not. We know of no closed form for this flight; the method of
§\ref{sec:method} computes it to any accuracy.

\subsection{What energy alone says}

The energy \xeqref{foundations}{eq:energy} with the potential energy
$V(\vect{x}) = \sum_i \kappa_i/\abs{\vect{x}-\vect{r}_i}$ is
\begin{equation}\label{eq:energy-two}
  W = \gamma mc^2 + V(\vect{x}) = \gamma_0 mc^2 + V(\vect{x}_0).
\end{equation}
Since $\gamma \ge 1$, the particle can only be where
\begin{equation}\label{eq:allowed}
  V(\vect{x}) \le W - mc^2 = T_0 + V(\vect{x}_0).
\end{equation}
At the launch, $\abs{\vect{x}_0 - \vect{r}_1} = \val{rOneZero}$ and
$\abs{\vect{x}_0-\vect{r}_2} = \val{rTwoZero}$, so $V(\vect{x}_0) = \val{VZero}$ and
$W - mc^2 = \val{Wbind}$: negative. As $V \to 0$ far away, \eqref{eq:allowed} confines the
particle to a bounded region around the two charges (Figure~\ref{fig:board}): it cannot escape,
and it must turn back. Energy says no more: where in that region it goes, and whether it ever
reaches the detector, needs the equations of motion.

\subsection{The method}\label{sec:method}

The state $y = (x, y, p_x, p_y)$ evolves by \eqref{eq:eom-two}. The classical Runge--Kutta
method \xeqref{foundations}{eq:rk4} advances it by a fixed step $h$, its right-hand side
evaluated four times per step. After each step, check whether the step's path met an event
boundary: the detector's box, a sphere of radius $R + a$ about either charge, the world's edge.
On the step $[t_n, t_n + h]$ the position is interpolated by the cubic polynomial that matches the
positions and velocities at both ends (\xsecref{foundations}{sec:numerics}),
\begin{equation}\label{eq:hermite}
  \vect{x}(t_n + uh) = h_{00}(u)\,\vect{x}_n + h_{10}(u)\,h\vect{v}_n + h_{01}(u)\,\vect{x}_{n+1} + h_{11}(u)\,h\vect{v}_{n+1},
\end{equation}
with $h_{00} = 2u^3 - 3u^2 + 1$, $h_{10} = u^3 - 2u^2 + u$, $h_{01} = -2u^3 + 3u^2$,
$h_{11} = u^3 - u^2$ and $\vect{v} = \vect{p}/(\gamma m)$, sampled at 64 points of the step;
bisection on $u$ then finds the first instant inside the box. The state at that instant comes
from one more step of the method, of length $t_\mathrm{e} - t_n$. The method's error is of order
$h^4$; repeating the computation with $h/2$ estimates it as a fifteenth of the difference
\xsecref{foundations}{sec:numerics}. The energy \eqref{eq:energy-two} gives an independent check
of every step.

\section{The reference solution by hand}\label{sec:hand}

\begin{procedure}
  \item \emph{The launch} \xeqref{foundations}{eq:launch}: $\gamma_0 = \val{gammaZero}$ and
    $\abs{\vect{p}_0} = \val{pZero}$, so $\vect{p}_0 = (0, \val{pZero})$ up to the rounding of the
    direction.
  \item \emph{The energy} \eqref{eq:energy-two}: $W = \val{W}$ and $W - mc^2 = \val{Wbind}$: the
    particle is bound \eqref{eq:allowed}.
  \item\label{st:integrate} \emph{The flight}: the method of §\ref{sec:method} with
    $h = \val{HTable}$, from $t = 0$. The states every five time units, with the energy's drift
    as a check of the arithmetic:
    \begin{center}\small\val{states}\end{center}
    The particle climbs to the upper left, slows to its smallest momentum near $(0.9, 12.3)$,
    where it almost touches the edge of the allowed region (Figure~\ref{fig:board}), and falls
    back towards $Q_2$.
  \item \emph{The entry}: with $h = \val{H}$ the step from $t_n = \val{tBefore}$, at
    $(\val{xBefore}, \val{yBefore})$, ends at $(\val{xAfter}, \val{yAfter})$, inside the detector,
    which the path entered through its \val{faceE} face $y = 4$. The interpolant
    \eqref{eq:hermite} crosses $y = 4$ at $u = \val{uE}$: $t_\mathrm{e} = \val{tE}$ and
    $x_\mathrm{e} = \val{xE}$, after \val{steps} steps.
  \item \emph{The momentum at the entry}: one step of length $t_\mathrm{e} - t_n$ gives
    $\vect{p}_\mathrm{e} = (\val{pxE}, \val{pyE})$. The energy alone gives
    $\gamma_\mathrm{e} = (W - V(\vect{x}_\mathrm{e}))/(mc^2) = \val{gammaE}$ and
    $\abs{\vect{p}_\mathrm{e}} = mc\sqrt{\gamma_\mathrm{e}^2-1} = \val{pEnergy}$, which the
    integration matches to $\val{pDiff}$.
  \item \emph{The error}: halving the step estimates the error of $t_\mathrm{e}$ as $\val{errT}$
    and of $x_\mathrm{e}$ as $\val{errX}$; along the whole flight the energy stays within
    $\val{wDev}$ of $W$. The step can be much coarser for a rough answer:
    \begin{center}\small\val{convergence}\end{center}
  \item \emph{The radiation}: the power \xeqref{foundations}{eq:larmor} with the acceleration
    \xeqref{first_bend}{eq:acceleration}, summed over the steps by Simpson's rule
    \xeqref{foundations}{eq:simpson}: $E_\mathrm{rad} = \val{Erad}$.
\end{procedure}

\begin{result}[The reference flight]
The particle enters the detector through its top face at $(\val{xE}, 4)$ at
$t_\mathrm{e} = \val{tE}$, with momentum $\vect{p}_\mathrm{e} = (\val{pxE}, \val{pyE})$: the
level is solved.
\end{result}

By hand, Step~\ref{st:integrate} is long: $\val{tableSteps}$ steps of four evaluations of
\eqref{eq:eom-two} each at $h = \val{HTable}$. The convergence table shows what a patient
calculator gains from fewer: at $h = 0.04$, $\val{roughSteps}$ steps, the arrival is still right
to $\val{roughErr}$.

\section{Comparison with the game}\label{sec:comparison}

\comparisontable

The method's own error is about $10^{-12}$ (Step~6), so the differences in the table are the
game's: below $10^{-10}$ for the reference flight, as its verified integration promises
(\xsecref{foundations}{sec:verification}). The largest, $\val{cmpDetThreeTDiff}$, belongs to the
longest flight, placement~3's of $\val{detThreeT}$ time units, over which the game's error grows;
the tolerance of $10^{-9}$, set from the game's accuracy before comparing, still holds it. The
radiated energy agrees to $\val{cmpRadDiff}$ (relative).

\section{Other solutions}\label{sec:others}

\begin{figure}[htbp]
\centering
\begin{tikzpicture}
\begin{groupplot}[group style={group size=3 by 2, horizontal sep=0.6cm, vertical sep=1.0cm},
  axis equal image, width=0.37\linewidth, xmin=-1, xmax=31, ymin=-1, ymax=21,
  xtick={0,10,20,30}, ytick={0,10,20}, tick label style={font=\scriptsize},
  title style={font=\small, yshift=-4pt}]
  \nextgroupplot[title={$Q = \num{1e6}$: \val{nplusone}}]
    \fill[green!50!black, fill opacity=0.2] (0,0) rectangle (3,4);
    \draw[black!40] (0,0) rectangle (30,20);
    \fill[blue!60!black] (15,12) circle[radius=0.5];
    \addplot[only marks, mark=*, mark size=1.1pt, red!70!black] table {map_plusone.dat};
    \fill (6,3) circle[radius=0.4];
  \nextgroupplot[title={$Q = \num{2e6}$: \val{nplustwo}}]
    \fill[green!50!black, fill opacity=0.2] (0,0) rectangle (3,4);
    \draw[black!40] (0,0) rectangle (30,20);
    \fill[blue!60!black] (15,12) circle[radius=0.5];
    \addplot[only marks, mark=*, mark size=1.1pt, red!70!black] table {map_plustwo.dat};
    \fill (6,3) circle[radius=0.4];
  \nextgroupplot[title={$Q = \num{4e6}$: \val{nplusfour}}]
    \fill[green!50!black, fill opacity=0.2] (0,0) rectangle (3,4);
    \draw[black!40] (0,0) rectangle (30,20);
    \fill[blue!60!black] (15,12) circle[radius=0.5];
    \addplot[only marks, mark=*, mark size=1.1pt, red!70!black] table {map_plusfour.dat};
    \fill (6,3) circle[radius=0.4];
  \nextgroupplot[title={$Q = \num{-1e6}$: \val{nminusone}}]
    \fill[green!50!black, fill opacity=0.2] (0,0) rectangle (3,4);
    \draw[black!40] (0,0) rectangle (30,20);
    \fill[blue!60!black] (15,12) circle[radius=0.5];
    \addplot[only marks, mark=*, mark size=1.1pt, blue!70!black] table {map_minusone.dat};
    \fill (6,3) circle[radius=0.4];
  \nextgroupplot[title={$Q = \num{-2e6}$: \val{nminustwo}}]
    \fill[green!50!black, fill opacity=0.2] (0,0) rectangle (3,4);
    \draw[black!40] (0,0) rectangle (30,20);
    \fill[blue!60!black] (15,12) circle[radius=0.5];
    \addplot[only marks, mark=*, mark size=1.1pt, blue!70!black] table {map_minustwo.dat};
    \fill (6,3) circle[radius=0.4];
  \nextgroupplot[title={$Q = \num{-4e6}$: \val{nminusfour}}]
    \fill[green!50!black, fill opacity=0.2] (0,0) rectangle (3,4);
    \draw[black!40] (0,0) rectangle (30,20);
    \fill[blue!60!black] (15,12) circle[radius=0.5];
    \addplot[only marks, mark=*, mark size=1.1pt, blue!70!black] table {map_minusfour.dat};
    \fill (6,3) circle[radius=0.4];
    \draw[red!70!black, thick] (2,0) circle[radius=0.9];
\end{groupplot}
\end{tikzpicture}
\caption{The game's verdicts: every node where a single added charge solves the level, for each
charge the player may place (repelling red, attracting blue), with the level's fixed charge $Q_1$
(large dot) and the launch (black). The ring marks the reference. These maps are the game's, not
derived here; three of their placements are checked below.}
\label{fig:map}
\end{figure}

Of the $\val{nFamily}$ single-charge placements, the game finds $\val{nSolving}$ that solve the
level (Figure~\ref{fig:map}). Most are strong attracting charges below and to the left of the
launch, which catch the particle as $Q_2$ does; a few are repelling charges above it, which push it
back down. The method of §\ref{sec:method} checks three of them independently (\code{placements.json};
each with its error estimate from halving the step):
\begin{enumerate}[label=\arabic*.]
  \item $Q = \val{detOneQ}$ at $\val{detOneNode}$: through the \val{detOneFace} face at
    $(\val{detOneX}, \val{detOneY})$, at $t = \val{detOneT}$ (error $\val{detOneErr}$).
  \item $Q = \val{detTwoQ}$ at $\val{detTwoNode}$, repelling: through the \val{detTwoFace} face at
    $(\val{detTwoX}, \val{detTwoY})$, at $t = \val{detTwoT}$ (error $\val{detTwoErr}$), pushed straight
    back before it climbs.
  \item $Q = \val{detThreeQ}$ at $\val{detThreeNode}$: through the \val{detThreeFace} face at
    $(\val{detThreeX}, \val{detThreeY})$, at $t = \val{detThreeT}$ (error $\val{detThreeErr}$), after
    several loops about the two attracting charges.
\end{enumerate}
With two or three charges the possibilities multiply; each is a flight of the same kind, and the
method above computes any of them.

\section{What the game computes}\label{sec:engine}

The game integrates \eqref{eq:eom-two} with its eighth-order Runge--Kutta method and step-size
control at the tolerances $10^{-10}$ and $10^{-12}$ (\physref{5}). Where this document takes
\val{steps} steps of a fourth-order method, the game's verified run takes
$\val{gameSteps}$: large steps where the force changes slowly, small ones near the charges. It
finds the entry into the detector on its method's interpolant with a search that cannot miss a
crossing (\physref{6}), and verifies the outcome by comparing its two runs (\physref{7}). The model
is this document's; only \ref{as:radiation} is an approximation, and the game reports its size
(\physref{8}).

\section{Further reading}\label{sec:reading}

\begin{itemize}
  \item \levelref{first_bend}: the flight about one charge, in closed form, which this one
    reduces to near either charge.
  \item \HairerI{§§II.1 and II.4}: the Runge--Kutta methods, their errors and step-size control;
    §II.6 for dense output.
  \item \Goldstein{Ch.~10}: the Hamilton--Jacobi equation and the separation of variables that
    solves the Newtonian problem of two fixed centres, and fails here.
\end{itemize}

\levelfileappendix

\end{document}
